\documentclass{article}
\usepackage{amsmath, amssymb}

\begin{document}

\section*{Pushing--Medium Gravity with Effective Metric (Static Formulation v1)}

This document summarizes the gravitational sector of the Pushing--Medium (PM)
model in the static limit, including the effective isotropic metric that
emerges from the PM refractive index field. Only mass density and mass currents
are included.

The PM gravitational field is encoded in two static fields:



\[
\phi(\mathbf r) \equiv \ln n(\mathbf r),
\qquad
\mathbf u(\mathbf r) \equiv \text{gravitational flow field}.
\]



The scalar field $\phi$ determines the refractive index $n$, while the flow
field $\mathbf u$ encodes moving-mass effects.

\section*{1. Gravitational Index Field}

The scalar field $\phi = \ln n$ satisfies a Poisson equation sourced by mass
density:



\[
\nabla^2 \phi(\mathbf r)
=
-4\pi G_{\mathrm{eff}}\,\rho_M(\mathbf r).
\]



The gravitational refractive index is:



\[
n(\mathbf r) = \exp(\phi(\mathbf r)).
\]



\subsection*{1.1 Point Mass}

For a point mass $M$ at the origin:



\[
\phi(r) = \frac{\mu_G M}{r},
\qquad
n(r) = 1 + \frac{\mu_G M}{r} + \mathcal{O}(r^{-2}).
\]



This reproduces the PM gravitational index used in bending calculations.

\subsection*{1.2 Continuous Mass Distribution}

For a general mass density $\rho_M(\mathbf r)$:



\[
\phi(\mathbf r)
=
\mu_G
\int
\frac{\rho_M(\mathbf r')}{|\mathbf r - \mathbf r'|}
\,d^3r'.
\]



\section*{2. Gravitational Flow Field}

The static flow field $\mathbf u(\mathbf r)$ is sourced by mass currents:



\[
\nabla^2 \mathbf u(\mathbf r)
=
- A_M\,\mathbf J_M(\mathbf r).
\]



This field encodes:

\begin{itemize}
    \item moving-mass lensing,
    \item Fizeau-like dragging,
    \item frame-dragging analogues.
\end{itemize}

\subsection*{2.1 Moving Point Mass}

For a point mass $M$ moving with velocity $\mathbf v$:



\[
\mathbf u(\mathbf r)
=
A_M M
\frac{\mathbf v}{|\mathbf r - \mathbf r_M|}.
\]



\section*{3. Effective Isotropic Metric}

In the weak-field limit, the PM refractive index $n = e^\phi$ induces an
effective isotropic metric of the form:



\[
ds^2
=
-(1 + 2\phi)\,c^2 dt^2
+
(1 - 2\phi)\,(dx^2 + dy^2 + dz^2).
\]



This metric reproduces the standard post-Newtonian expansion of general
relativity in isotropic coordinates.

\subsection*{3.1 PPN Parameters}

From the effective metric:



\[
g_{tt} = -(1 + 2\phi),
\qquad
g_{rr} = (1 - 2\phi),
\]



the PPN parameters are:



\[
\gamma = 1,
\qquad
\beta = 1.
\]



These match the general-relativistic values in the weak-field regime.

\section*{4. Ray and Particle Propagation}

In static gravitational fields, rays follow:



\[
\frac{d\mathbf r}{ds}
= c\,\hat{\mathbf k}
+ k_{\mathrm{Fizeau}}\,\mathbf u(\mathbf r),
\]





\[
\frac{d\hat{\mathbf k}}{ds}
= \nabla_\perp \phi(\mathbf r)
+ \mathcal{F}\big(\hat{\mathbf k}, \mathbf u, \nabla \mathbf u\big).
\]



The bending angle in the weak/semistrong regime satisfies:



\[
\alpha(b) \propto \frac{1}{b},
\]



consistent with numerical tests and with the PM gravitational index field.

\section*{5. Summary}

The PM gravitational sector consists of:

\begin{itemize}
    \item a scalar index field sourced by mass density,
    \item a flow field sourced by mass currents,
    \item an effective isotropic metric with $\gamma = \beta = 1$,
    \item weak/semistrong bending with $\alpha \propto 1/b$,
    \item moving-mass and Fizeau-like effects.
\end{itemize}

This formulation matches all current numerical tests in the weak and
semistrong regimes and provides the gravitational backbone of the PM model.

\end{document}

